\section{Introduction}
\IEEEPARstart{D}{eep} Neural Networks (DNNs) have established foundational breakthroughs across core computer vision tasks, ranging from large-scale visual recognition \cite{11, 24} to dense spatial reasoning and object detection \cite{14, 53}. This trajectory was catalyzed by Convolutional Neural Networks (CNNs) \cite{23}, which embedded translation equivariance and local spatial hierarchies into learnable architectures. In recent years, the computer vision landscape has been profoundly shaped by the advent of Vision Transformers (ViTs) \cite{9}, masked image modeling \cite{46}, and self-supervised foundation backbones \cite{1}. However, this contemporary trajectory has predominantly reinforced a ``scale-driven'' paradigm, wherein marginal performance increments are achieved by training networks with tens or hundreds of millions---and increasingly billions---of parameters over colossal compute clusters and proprietary web-scale datasets \cite{11, 42}.

In parallel, deep learning has revolutionized medical image analysis \cite{10, 27}, demonstrating specialist-level efficacy in complex diagnostic tasks \cite{10}. Within neuro-oncology, non-invasive Magnetic Resonance Imaging (MRI) serves as the primary clinical imaging modality for the detection, differential diagnosis, and histologic grading of intracranial neoplasms, including gliomas, meningiomas, and pituitary adenomas \cite{2, 34}. Deep learning architectures applied to multi-sequence brain MRI have automated diagnostic workflows that historically required exhaustive manual radiological review \cite{2, 34, 38}. Standard clinical diagnostic paradigms rely heavily on overparameterized deep backbones---such as deep ResNets \cite{14}, VGG networks \cite{39}, and compound-scaled EfficientNets \cite{40}---to identify subtle textural heterogeneity, necrotic cores, and invasive margins across disparate axial, coronal, and sagittal planes \cite{36, 44}.

Despite these empirical successes, a critical translational divide persists between benchtop benchmark accuracy and clinical, point-of-care utility \cite{20, 26}. In healthcare environments, real-world deployment rarely occurs in centralized, high-throughput cloud data centers. Routine transmission of high-resolution volumetric neuroimaging to external cloud infrastructures encounters severe regulatory barriers, notably patient confidentiality mandates and data protection governance such as HIPAA in the United States and GDPR in the European Union \cite{10, 26}. Furthermore, cloud dependency introduces vulnerability to variable network bandwidth, transmission latency, and service outages---risks that are unacceptable during acute clinical scenarios, emergency neuro-triage, or intra-operative surgical navigation \cite{20, 33}. Consequently, bringing intelligence directly to the clinical edge---operating locally on hospital workstations, embedded scanner consoles, and mobile intra-operative carts---has emerged as an urgent imperative \cite{12, 25, 26}.

However, edge deployment imposes uncompromising physical constraints: strict inference latency windows, severely constrained Random Access Memory (RAM), low thermal design power (TDP), and limited on-chip arithmetic hardware \cite{12, 33, 41}. To accommodate mobile and decentralized environments, the vision community introduced lightweight architectural designs, including depthwise separable convolutions in MobileNets \cite{15, 37}, channel shuffling in ShuffleNet \cite{32}, and parameter squeezing in SqueezeNet \cite{17}. Yet, the overwhelming majority of these ``efficient'' models predominantly inhabit a mobile-scale budget of 2 to 20 million parameters. In contrast, the sub-1M parameter regime ($<1,000,000$ parameters) remains fundamentally underexplored in medical imaging, particularly under the practical constraints of small-to-medium institutional datasets \cite{3, 21}. Crucially, as observed in recent foundational studies on compact architectures \cite{21}, naively downscaling standard architectures via uniform width or depth multipliers leads to erratic representational capacity, gradient instability, and disproportionate performance collapse \cite{21, 24}.

This investigation is motivated by the absence of systematic benchmarks, architectural insights, and deployment guidelines for deep learning models operating under a strict sub-1M parameter ceiling in neuro-oncological imaging. We posit that as the parameter budget is aggressively constrained, architectural topology---specifically how parameters are structurally allocated to capture multi-scale context without parameter redundancy---matters far more than sheer parameter volume \cite{21}. While multi-path atrous (dilated) depthwise convolutions have emerged as a premier mechanism for expanding effective receptive fields without parameter inflation \cite{5, 21}, their behavior, domain robustness, and failure mechanics under clinical distribution shifts have never been systematically characterized.

To address these foundational questions, this work conducts a systematic empirical study of sub-1M parameter architectures across two clinically grounded neuro-oncology benchmarks: (i) an in-domain, 4-class multi-tumor diagnostic classification task (the Nickparvar benchmark \cite{34}), and (ii) a challenging out-of-domain cross-dataset glioma grading task (BraTS 2020 \cite{2, 28}) exhibiting severe class imbalance and multi-sequence distribution shifts. Our empirical findings uncover two critical phenomena:
\begin{enumerate}
\item \textbf{The In-Domain Saturation Phenomenon}: On saturated diagnostic benchmarks, architecturally optimized sub-1M models achieve diagnostic discrimination statistically comparable to models an order of magnitude larger. Specifically, a compact multi-path atrous architecture (LiteSpeed-Net, LS-Net \cite{21}) achieves 98.29\% accuracy at only 0.36 GFLOPs, operating six times faster ($6.0\times$ speedup, 11.3 ms CPU latency, 88.5 FPS) than a fine-tuned ResNet-50 (23.52M parameters, 8.17 GFLOPs, 67.89 ms latency) while forfeiting less than one percentage point of accuracy ($<0.85$ pp).
\item \textbf{The Multi-Scale Cross-Dataset Vulnerability}: When exposed to challenging cross-dataset clinical grading, the saturation breaks, exposing an acute vulnerability in compact multi-path atrous convolutions. Because compact networks rely on parallel dilated branches to capture receptive field coverage without learned parameter depth, their predictions are disproportionately brittle under input scale perturbations. Diagnostic errors are highly concentrated in low-confidence predictions located near ambiguous lesion boundaries.
\end{enumerate}

Crucially, in resource-constrained clinical settings, scaling up parameter depth or deploying heavy ensemble backbones to resolve this residual error is computationally impermissible \cite{21, 41}. Instead, we argue that the residual accuracy gap can be recovered far more efficiently at inference time by exploiting the model's own architectural inductive biases. To this end, we introduce \textbf{Multi-Scale Consistency Gated Test-Time Augmentation (MCG-TTA)}. MCG-TTA is an architecture-informed, uncertainty-aware inference mechanism designed to stabilize scale-sensitive compact representations with \textbf{zero additional learned parameters} (retaining a 0.91M parameter footprint). Rather than executing costly, unconditional multi-view augmentation across all inputs \cite{4, 47}, MCG-TTA monitors the model's prediction confidence. Samples exhibiting high predictive uncertainty ($p_{\max} < \tau$) are selectively probed with scale perturbations ($\pm 10\%$), directly mirroring the receptive field offsets of the network's internal atrous dilation branches. Multi-view aggregation is selectively triggered only when scale disagreement is detected.

Extensive evaluations confirm that MCG-TTA recovers up to $+1.8$ percentage points in slice-level F1-score on out-of-domain BraTS grading while operating at only $\sim 1.9\times$ baseline latency, while preserving 100\% of detected low-grade glioma (LGG) clinical patient diagnoses (20/24). These results demonstrate that uncertainty-aware test-time calibration offers a superior, hardware-friendly alternative to model scaling for clinical edge artificial intelligence.

In summary, this work provides the following key contributions:
\begin{itemize}
\item \textbf{Comprehensive Sub-1M Medical Imaging Benchmark}: We provide the first systematic benchmarking of deep learning architectures strictly constrained to $<1\text{M}$ parameters across both in-domain multi-tumor classification (Nickparvar) and out-of-domain clinical glioma grading (BraTS 2020), benchmarking against mobile-scale and heavy overparameterized baselines (ResNet-50, VGG-16, EfficientNet-B0).
\item \textbf{Characterization of Saturation and Failure Modes}: We empirically demonstrate that sub-1M models achieve near-lossless diagnostic parity ($<1$ pp difference) with $6\times$ speedup on saturated benchmarks, and formally diagnose the failure mechanism of compact multi-path atrous convolutions under cross-dataset scale perturbations.
\item \textbf{Architecture-Informed MCG-TTA}: We propose Multi-Scale Consistency Gated Test-Time Augmentation, an uncertainty-gated inference strategy that mirrors internal atrous dilation scales to recover residual accuracy selectively (+1.8 pp slice F1 at $1.9\times$ latency) with 0 additional learned parameters.
\item \textbf{Rigorous Edge Profiling and Clinical Safety Validation}: We profile floating-point operations (FLOPs), multiply-accumulate operations (MACs), real single-batch CPU latency, and throughput (FPS), alongside patient-level clinical sensitivity, verifying that compact uncertainty-aware inference delivers robust real-time diagnosis on resource-constrained hardware.
\end{itemize}
